\clearpage
\oldpage[IV-3]{249}
\phantomsection
\addcontentsline{toc}{section}{Bibliografía histórica (continuación)}
\section*{Bibliografía histórica (continuación)}

\begin{description}[style=multiline,leftmargin=3.5em,labelwidth=2.8em]
  \item[{[39]}] J.-P. Serre, \emph{Corps locaux}, París (Hermann), 1962.
  \item[{[40]}] J.-P. Serre, \emph{Groupes proalgébriques},
  \emph{Publ. Inst. Hautes Études Sci.}, n.º~7 (1960).
\end{description}

\clearpage
\oldpage[IV-3]{250}
\phantomsection
\addcontentsline{toc}{section}{Índice de notaciones}
\section*{Índice de notaciones}

\begingroup
\small
\setlength{\parindent}{0pt}
\setlength{\parskip}{2pt}

$\mathfrak{L}(X)$, $\mathfrak{O}c(X)$, $\mathfrak{G}c(X)$,
$\mathfrak{F}gc(X)$ ($X$ preesquema): 8.3.9.

$Q(\mathcal F)$ ($\mathcal F$ $\mathcal O_X$-Módulo): 8.5.10.

$\operatorname{Spr}(Y)$, $\operatorname{Spro}(Y)$,
$\operatorname{Sprf}(Y)$ ($Y$ preesquema): 8.6.1.

$\operatorname{Pro}(C)$ ($C$ categoría): 8.13.3.

$X_S$, $\mathcal F_S$, $u_S$, $s_S$, $Z_S$, $h_S$
($X$ $S$-preesquema, $s\in S$, $\mathcal F$ $\mathcal O_X$-Módulo,
$u$ homomorfismo de $\mathcal O_X$-Módulos, $s$ sección de un
$\mathcal O_X$-Módulo, $Z$ parte de $X$, $h:X\to Y$ $S$-morfismo): 9.4.1.

$Qc(X)$, $Sqc(X)$, $\mathfrak O(X)$, $\mathfrak G(X)$,
$\mathfrak U(X)$, $\mathfrak C(X)$ ($X$ espacio topológico): 10.1.1.

$\operatorname{prof}^{*}_{x}(\mathcal F)$,
$\operatorname{prof}_{*}(\mathcal F)$,
$\operatorname{prof}^{*}(\mathcal F)$
($\mathcal F$ $\mathcal O_X$-Módulo): 10.8.1.

$S(X)$ ($X$ preesquema de Jacobson): 10.9.1.

$S(f)$ ($f$ morfismo de preesquemas de Jacobson): 10.9.2.

$\operatorname{Spm}(A)$, $D_m(a)$ ($A$ anillo de Jacobson, $a\in A$): 10.9.3.

$\operatorname{Tor.dim}_{B}(M)$ ($M$ $B$-módulo): 12.3.2.
\endgroup

\clearpage
\oldpage[IV-3]{251}
\phantomsection
\addcontentsline{toc}{section}{Índice terminológico}
\section*{Índice terminológico}

\begingroup
\small
\setlength{\parindent}{0pt}
\setlength{\parskip}{2pt}
\newcommand{\EGAIVthreeTerm}[2]{\textbf{#1}: #2.\par}

\EGAIVthreeTerm{Criterio de planitud por las fibras}{11.3.10}
\EGAIVthreeTerm{Espacio algebraico sobre $k$, espacio $k$-algebraico,
espacio prealgebraico sobre $k$, espacio $k$-prealgebraico}{10.10.2}
\EGAIVthreeTerm{Espacio de Jacobson}{10.3.1}
\EGAIVthreeTerm{Espacio $k$-prealgebraico reducido}{10.10.4}
\EGAIVthreeTerm{Familia de morfismos esquemáticamente dominante}{11.10.2}
\EGAIVthreeTerm{Familia de morfismos universalmente esquemáticamente dominante}{11.10.7}
\EGAIVthreeTerm{Familia de subpreesquemas esquemáticamente densa}{11.10.2}
\EGAIVthreeTerm{Familia separante de homomorfismos de haces de módulos}{11.9.1}
\EGAIVthreeTerm{Familia separante de homomorfismos de módulos}{11.9.4}
\EGAIVthreeTerm{Familia universalmente separante de homomorfismos de haces de módulos}{11.9.14}
\EGAIVthreeTerm{Geométricamente isomorfos (preesquemas sobre $k$)}{9.1.4}
\EGAIVthreeTerm{Homomorfismo de Módulos universalmente inyectivo}{11.9.14 y 11.9.18}
\EGAIVthreeTerm{Lema de Chow para los morfismos de presentación finita}{8.10.5.1}
\EGAIVthreeTerm{Teorema principal de Zariski para los morfismos de presentación finita}{8.12.6}
\EGAIVthreeTerm{Morfismo equidimensional en un punto, morfismo equidimensional}{13.2.2 y 13.3.2}
\EGAIVthreeTerm{Morfismo propio en un punto}{15.7.1}
\EGAIVthreeTerm{Morfismo pseudofinito}{8.12.3}
\EGAIVthreeTerm{Morfismo submersivo, morfismo universalmente submersivo}{15.7.8}
\EGAIVthreeTerm{Morfismo universalmente abierto en un punto}{14.3.3}
\EGAIVthreeTerm{Abierto ultraafín}{10.9.5}
\EGAIVthreeTerm{Parte finita saturada de un preesquema algebraico sobre $k$}{9.8.9}
\EGAIVthreeTerm{Parte propia en un punto (de un $S$-preesquema)}{15.7.1}
\EGAIVthreeTerm{Parte cuasiconstruible, parte localmente cuasiconstruible de un espacio topológico}{10.1.1}
\EGAIVthreeTerm{Parte muy densa de un espacio topológico}{10.1.3}
\EGAIVthreeTerm{Polinomio geométricamente irreducible}{9.7.4}
\EGAIVthreeTerm{Preesquema de Jacobson}{10.4.1}
\EGAIVthreeTerm{Preesquema equidimensional sobre otro en un punto, equidimensional sobre otro}{13.2.2 y 13.3.2}
\EGAIVthreeTerm{Preesquema esencialmente afín sobre otro}{8.13.4}
\EGAIVthreeTerm{Principio de extensión finita}{9.1.1}
\EGAIVthreeTerm{Profundidad rectificada}{10.8.1}
\EGAIVthreeTerm{Proobjeto, morfismo de proobjetos}{8.13.3}
\EGAIVthreeTerm{Proobjeto esencialmente afín}{8.13.4, 8.14.1}
\EGAIVthreeTerm{Propiedad construible, propiedad ind-construible}{9.2.1 y 9.2.2, (i) y (ii)}
\EGAIVthreeTerm{Cuasihomeomorfismo de espacios topológicos}{10.2.2}
\EGAIVthreeTerm{Cuasiisomorfismo de espacios anillados}{10.2.8}
\EGAIVthreeTerm{Saturado de una parte finita de un preesquema algebraico sobre $k$}{9.8.9}
\EGAIVthreeTerm{Subconjunto de un $S$-preesquema submersivo, universalmente submersivo sobre su imagen}{15.7.8}
\EGAIVthreeTerm{Esqueleto primario de un Módulo, esqueleto virtual}{9.8.9}
\EGAIVthreeTerm{Espectro maximal de un anillo de Jacobson}{10.9.3}
\EGAIVthreeTerm{Tipo primario de un Módulo, tipo primario geométrico de un Módulo}{9.8.9}
\EGAIVthreeTerm{Ultrapreesquema, morfismo de ultrapreesquemas}{10.9.5}
\EGAIVthreeTerm{$R$-ultrapreesquema}{10.10.2}
\endgroup

\clearpage
\oldpage[IV-3]{253}
\phantomsection
\addcontentsline{toc}{section}{Tabla de materias original}
\section*{Tabla de materias original}

\begingroup
\small
\newcommand{\EGAIVthreeContents}[3]{\textbf{#1}\quad #2\dotfill #3\par}

\EGAIVthreeContents{Capítulo IV}{Estudio local de los esquemas y de los morfismos de esquemas (continuación)}{5}
\EGAIVthreeContents{§8}{Límites proyectivos de preesquemas}{5}
\EGAIVthreeContents{8.1}{Introducción}{5}
\EGAIVthreeContents{8.2}{Límites proyectivos de preesquemas}{7}
\EGAIVthreeContents{8.3}{Partes construibles en un límite proyectivo de preesquemas}{12}
\EGAIVthreeContents{8.4}{Criterios de irreducibilidad y conexión para los límites proyectivos de preesquemas}{17}
\EGAIVthreeContents{8.5}{Módulos de presentación finita sobre un límite proyectivo de preesquemas}{19}
\EGAIVthreeContents{8.6}{Subpreesquemas de presentación finita de un límite proyectivo de preesquemas}{25}
\EGAIVthreeContents{8.7}{Criterios para que un límite proyectivo de preesquemas sea reducido (resp. íntegro)}{27}
\EGAIVthreeContents{8.8}{Preesquemas de presentación finita sobre un límite proyectivo de preesquemas}{28}
\EGAIVthreeContents{8.9}{Primeras aplicaciones a la eliminación de hipótesis noetherianas}{34}
\EGAIVthreeContents{8.10}{Propiedades de permanencia de los morfismos por paso al límite proyectivo}{36}
\EGAIVthreeContents{8.11}{Aplicación a los morfismos cuasifinitos}{41}
\EGAIVthreeContents{8.12}{Nueva demostración y generalización del teorema principal de Zariski}{43}
\EGAIVthreeContents{8.13}{Traducción en términos de proobjetos}{49}
\EGAIVthreeContents{8.14}{Caracterización de un preesquema localmente de presentación finita sobre otro mediante el funtor que representa}{52}

\EGAIVthreeContents{§9}{Propiedades construibles}{54}
\EGAIVthreeContents{9.1}{El principio de extensión finita}{54}
\EGAIVthreeContents{9.2}{Propiedades construibles e ind-construibles}{56}
\EGAIVthreeContents{9.3}{Propiedades construibles de morfismos de preesquemas algebraicos}{60}

\oldpage[IV-3]{254}
\EGAIVthreeContents{9.4}{Construibilidad de ciertas propiedades de los Módulos}{62}
\EGAIVthreeContents{9.5}{Construibilidad de propiedades topológicas}{67}
\EGAIVthreeContents{9.6}{Construibilidad de ciertas propiedades de los morfismos}{71}
\EGAIVthreeContents{9.7}{Construibilidad de la separabilidad, la irreducibilidad geométrica y la conexión geométrica}{76}
\EGAIVthreeContents{9.8}{Descomposición primaria en un entorno de una fibra genérica}{83}
\EGAIVthreeContents{9.9}{Construibilidad de las propiedades locales de las fibras}{88}

\EGAIVthreeContents{§10}{Preesquemas de Jacobson}{95}
\EGAIVthreeContents{10.1}{Partes muy densas de un espacio topológico}{95}
\EGAIVthreeContents{10.2}{Cuasihomeomorfismos}{97}
\EGAIVthreeContents{10.3}{Espacios de Jacobson}{101}
\EGAIVthreeContents{10.4}{Preesquemas de Jacobson y anillos de Jacobson}{101}
\EGAIVthreeContents{10.5}{Preesquemas de Jacobson noetherianos}{104}
\EGAIVthreeContents{10.6}{Dimensión en los preesquemas de Jacobson}{107}
\EGAIVthreeContents{10.7}{Ejemplos y contraejemplos}{109}
\EGAIVthreeContents{10.8}{Profundidad rectificada}{110}
\EGAIVthreeContents{10.9}{Espectros maximales y ultrapreesquemas}{112}
\EGAIVthreeContents{10.10}{Espacios algebraicos de Serre}{114}

\EGAIVthreeContents{§11}{Propiedades topológicas de los morfismos planos de presentación finita. Criterios de planitud}{116}
\EGAIVthreeContents{11.1}{Lugares de planitud (caso noetheriano)}{117}
\EGAIVthreeContents{11.2}{Planitud de un límite proyectivo de preesquemas}{119}
\EGAIVthreeContents{11.3}{Aplicación a la eliminación de hipótesis noetherianas}{132}
\EGAIVthreeContents{11.4}{Descenso de la planitud por morfismos cualesquiera: caso de un preesquema base artiniano}{143}
\EGAIVthreeContents{11.5}{Descenso de la planitud por morfismos cualesquiera: caso general}{150}
\EGAIVthreeContents{11.6}{Descenso de la planitud por morfismos cualesquiera: caso de un preesquema base unirrama}{154}
\EGAIVthreeContents{11.7}{Contraejemplos}{157}
\EGAIVthreeContents{11.8}{Un criterio valuativo de planitud}{159}
\EGAIVthreeContents{11.9}{Familias separantes y universalmente separantes de homomorfismos de haces de módulos}{160}
\EGAIVthreeContents{11.10}{Familias esquemáticamente dominantes de morfismos y familias esquemáticamente densas de subpreesquemas}{170}

\EGAIVthreeContents{§12}{Estudio de las fibras de los morfismos planos de presentación finita}{173}
\EGAIVthreeContents{12.0}{Introducción}{173}
\EGAIVthreeContents{12.1}{Propiedades locales de las fibras de un morfismo plano localmente de presentación finita}{174}

\oldpage[IV-3]{255}
\EGAIVthreeContents{12.2}{Propiedades locales y globales de las fibras de un morfismo propio, plano y de presentación finita}{179}
\EGAIVthreeContents{12.3}{Propiedades cohomológicas locales de las fibras de un morfismo plano y localmente de presentación finita}{183}

\EGAIVthreeContents{§13}{Morfismos equidimensionales}{187}
\EGAIVthreeContents{13.1}{Teorema de semicontinuidad de Chevalley}{188}
\EGAIVthreeContents{13.2}{Morfismos equidimensionales: caso de morfismos dominantes de preesquemas irreducibles}{190}
\EGAIVthreeContents{13.3}{Morfismos equidimensionales: caso general}{194}

\EGAIVthreeContents{§14}{Morfismos universalmente abiertos}{199}
\EGAIVthreeContents{14.1}{Morfismos abiertos}{200}
\EGAIVthreeContents{14.2}{Morfismos abiertos y fórmula de las dimensiones}{202}
\EGAIVthreeContents{14.3}{Morfismos universalmente abiertos}{204}
\EGAIVthreeContents{14.4}{Criterio de Chevalley para los morfismos universalmente abiertos}{209}
\EGAIVthreeContents{14.5}{Morfismos universalmente abiertos y cuasisecciones}{216}

\EGAIVthreeContents{§15}{Estudio de las fibras de un morfismo universalmente abierto}{223}
\EGAIVthreeContents{15.1}{Multiplicidades de las fibras de un morfismo universalmente abierto}{223}
\EGAIVthreeContents{15.2}{Planitud de los morfismos universalmente abiertos con fibras geométricamente reducidas}{226}
\EGAIVthreeContents{15.3}{Aplicación: criterios de reducción e irreducibilidad}{228}
\EGAIVthreeContents{15.4}{Complementos sobre los morfismos de Cohen--Macaulay}{229}
\EGAIVthreeContents{15.5}{Rango separable de las fibras de un morfismo cuasifinito y universalmente abierto. Aplicación a las componentes conexas geométricas de las fibras de un morfismo propio}{231}
\EGAIVthreeContents{15.6}{Componentes conexas de las fibras a lo largo de una sección}{236}
\EGAIVthreeContents{15.7}{Apéndice: criterios valuativos de propiedad local}{242}
\EGAIVthreeContents{}{Bibliografía (continuación)}{249}
\EGAIVthreeContents{}{Índice de notaciones}{250}
\EGAIVthreeContents{}{Índice terminológico}{251}

\vspace{2em}
\begin{flushright}
Manuscrito recibido el 15 de noviembre de 1964.
\end{flushright}
\endgroup

\clearpage
\phantomsection
\addcontentsline{toc}{section}{Bibliografía consolidada de la edición}
\section*{Bibliografía consolidada de la edición}
\nocite{*}
\bibliography{the}
\bibliographystyle{amsalpha}
